% ============================================================
% ПРИЛОЖЕНИЕ «КОД ЗАКАЗА УСТРОЙСТВА»
% ============================================================

% Извлекаем суффикс (тип защиты) из \devicetype
\StrCount{\devicetype}{-}[\numdash]
\StrBehind[\numdash]{\devicetype}{-}[\suffix]

% Заголовок приложения
{\color{unimain}\section[(обязательное) Код заказа устройства «\devicetype»]{(обязательное)\\Код заказа устройства «\devicetype»}\label{app:order}}

% Макрос для вставки рисунка с двухстрочной подписью
\newcommand{\orderfig}[5][]{%
  \begin{figure}[H]
    \centering
    \vspace*{1cm}% отступ сверху
    \includegraphics[#1]{#2}
    \caption[#3]{#3\\#4}\label{#5}
  \end{figure}%
  \medskip
}

% Общее примечание для всех рисунков (одинаковый шрифт)
\newcommand{\note}{Примечание: типоисполнение «А» -- \suffix}

% Рисунки с кодами заказа
\orderfig[width=0.9\textwidth,height=0.9\textheight,keepaspectratio]%
  {\genpath/Приложение. Код заказа/Unit-M510.pdf}%
  {Код заказа устройства «ЮНИТ-М510-\suffix»}%
  {\note}{appcard:img1}

\orderfig[width=1\textwidth,height=1\textheight,keepaspectratio]%
  {\genpath/Приложение. Код заказа/Unit-M514.pdf}%
  {Код заказа устройства «ЮНИТ-М514-\suffix»}%
  {\note}{appcard:img2}

\orderfig[width=1\textwidth,height=1\textheight,keepaspectratio]%
  {\genpath/Приложение. Код заказа/Unit-M519.pdf}%
  {Код заказа устройства «ЮНИТ-М519-\suffix»}%
  {\note}{appcard:img3}

\clearpage